% ==============================================================================
% 04_criteria1_plos.tex — Criterion 1: Expected Learning Outcomes
% Score: 2
% ==============================================================================

\criterion{1}{ผลการเรียนรู้ที่คาดหวัง}{Expected Learning Outcomes}

% ------------------------------------------------------------------------------
\criterionsub{1.1 The programme to show that the expected learning outcomes are appropriately formulated in accordance with an established learning taxonomy, are aligned to the vision and mission of the university, and are known to all stakeholders.}

หลักสูตรวิศวกรรมศาสตรมหาบัณฑิต สาขาวิชาวิศวกรรมคอมพิวเตอร์และปัญญาประดิษฐ์ (หลักสูตรใหม่ พ.ศ.~2568) ของมหาวิทยาลัยราชภัฏอุตรดิตถ์ ได้ดำเนินการกำหนดและพัฒนาผลการเรียนรู้ระดับหลักสูตร (Programme Learning Outcomes: PLOs) ภายใต้กรอบมาตรฐานคุณภาพการศึกษาระดับบัณฑิตศึกษาและเกณฑ์การประเมินคุณภาพหลักสูตร โดยนำความเห็นของผู้มีส่วนได้ส่วนเสียเข้ามาประกอบการตัดสินใจในการยกร่างก่อนนำสู่การอนุมัติจากสภามหาวิทยาลัย (มติครั้งที่ 4/2568 ลงวันที่ 4 เมษายน พ.ศ.~2568) เพื่อเตรียมการเปิดรับนักศึกษาในปีการศึกษา 2568 (ในปีการศึกษา 2568 มีนักศึกษาเข้าศึกษาจริง 4 คน ในภาคการศึกษาที่ 2)

ในการพัฒนาหลักสูตรและ PLO มีการพิจารณาผ่านกลไกกำกับคุณภาพในระดับคณะและมหาวิทยาลัยตามลำดับ ตั้งแต่การให้ความเห็นชอบในระดับคณะ การกลั่นกรองโดยคณะกรรมการบัณฑิตศึกษาประจำมหาวิทยาลัย และการพิจารณาโดยสภาวิชาการ ก่อนเสนอสภามหาวิทยาลัยเพื่ออนุมัติให้เปิดดำเนินการ ซึ่งทำให้กรอบ PLO ที่ระบุในมคอ.2 มีความชัดเจน สอดคล้องกับนโยบายมหาวิทยาลัย และตรวจสอบย้อนกลับได้จากเอกสารการประชุมที่เกี่ยวข้อง (อ้างอิง AUNQA-1-2 และ AUNQA-1-4)

การเก็บรวบรวมข้อมูลความเห็นของผู้มีส่วนได้ส่วนเสีย แยกเป็นกลุ่มใหญ่ พร้อมวิธีการที่ใช้ ดังนี้
\begin{enumerate}
  \item \textbf{ผู้ใช้บัณฑิต/สถานประกอบการ} — เก็บข้อมูลจากผู้ประกอบการด้านดิจิทัลและปัญญาประดิษฐ์ในภาคเหนือตอนล่าง ผ่านแบบสอบถามและการสนทนากลุ่มแบบมีโครงสร้าง (focus group) ในระหว่างการยกร่างหลักสูตร
  \item \textbf{ศิษย์เก่า} — เก็บข้อมูลจากศิษย์เก่าหลักสูตรวิศวกรรมคอมพิวเตอร์ระดับปริญญาตรีที่มีแนวโน้มจะเป็นผู้ใช้งานบัณฑิตระดับบัณฑิตศึกษา ผ่านการสนทนากลุ่ม ร่วมกับการทบทวนแนวโน้มอุตสาหกรรมจากเอกสารนโยบายด้านกำลังคนด้านปัญญาประดิษฐ์ของภาครัฐ
  \item \textbf{องค์กรปกครองส่วนท้องถิ่นและชุมชน} — รับฟังความคาดหวังต่อการพัฒนากำลังคนด้านวิศวกรรมและเทคโนโลยีเพื่อพัฒนาท้องถิ่น ผ่านการหารือ/การมีส่วนร่วมของหน่วยงานที่เกี่ยวข้องในพื้นที่ (อปท.\ และเครือข่ายพันธมิตร) ตามบริบทการพัฒนาหลักสูตร
  \item \textbf{นักศึกษาและผู้สมัครเข้าศึกษา} — หลักสูตรอยู่ระหว่างเริ่มเปิดรับในรอบแรก จึงเน้นการเผยแพร่ข้อมูลหลักสูตรและคำอธิบาย PLO ผ่านเอกสารรับสมัคร การประชาสัมพันธ์ และแผนการปฐมนิเทศนักศึกษาเมื่อเริ่มภาคการศึกษา เพื่อให้ผู้เข้าศึกษาเข้าใจผลการเรียนรู้ที่คาดหวัง
  \item \textbf{อาจารย์ผู้รับผิดชอบหลักสูตรและอาจารย์ผู้เกี่ยวข้อง} — เก็บข้อมูลผ่านการประชุมคณะกรรมการพัฒนาหลักสูตร การทบทวนร่าง PLO/PEO และการอ้างอิงกรอบ AUN-QA version~4.0 ร่วมกับมาตรฐานคุณวุฒิระดับบัณฑิตศึกษา (TQF)
\end{enumerate}

หลักสูตรได้เผยแพร่รายละเอียดหลักสูตรตาม มคอ.2 ฉบับที่ผ่านการอนุมัติจากสภามหาวิทยาลัย แก่กลุ่มเป้าหมายผ่านช่องทางต่าง ๆ ได้แก่ \textbf{เว็บไซต์หลักสูตร} (\href{https://syweb.kidcbc.work}{syweb.kidcbc.work}) เว็บไซต์มหาวิทยาลัยราชภัฏอุตรดิตถ์ เว็บไซต์คณะเทคโนโลยีอุตสาหกรรม มหาวิทยาลัยราชภัฏอุตรดิตถ์ \textbf{เพจ Facebook ของหลักสูตร} ``\href{https://www.facebook.com/profile.php?id=61577236346855}{วิศวกรรมคอมพิวเตอร์และปัญญาประดิษฐ์}'' เอกสารแนะนำหลักสูตรสำหรับผู้สมัคร และช่องทางประชาสัมพันธ์ของมหาวิทยาลัย/คณะ เพื่อให้ผู้มีส่วนได้ส่วนเสียเข้าถึงข้อมูลหลักสูตรและผลการเรียนรู้ที่คาดหวังได้อย่างสม่ำเสมอ

การกำหนดผลการเรียนรู้ของหลักสูตรจัดทำให้สอดคล้องกับกรอบคุณวุฒิแห่งชาติ (มคอ.) สำหรับระดับบัณฑิตศึกษา ครอบคลุมห้าดาน ได้แก่ (1)~ด้านคุณธรรมจริยธรรม (2)~ด้านความรู้ (3)~ด้านทักษะทางปัญญา (4)~ด้านทักษะความสัมพันธ์และความรับผิดชอบ และ (5)~ด้านทักษะการวิเคราะห์เชิงตัวเลข การสื่อสาร และการใช้เทคโนโลยีสารสนเทศ โดยข้อความ PLO เมื่อพิจารณาร่วมกับโครงสร้างรายวิชา กระบวนการจัดการเรียนการสอน และแผนการประเมินผล แสดงให้เห็นการกระจายการเรียนรู้ในทุกด้านตามเจตนาของกรอบมคอ.

จากนั้นหลักสูตรนำความต้องการที่สำคัญจากการรับฟังผู้มีส่วนได้ส่วนเสียไปวิเคราะห์และจัดลำดับความสำคัญ และเชื่อมโยงกับ \textbf{วิสัยทัศน์ พันธกิจ อัตลักษณ์ และปรัชญาการจัดการศึกษา} ของมหาวิทยาลัย คณะที่ร่วมรับผิดชอบหลักสูตร และบริบทของสาขาวิชาวิศวกรรมคอมพิวเตอร์และปัญญาประดิษฐ์ จากนั้นกำหนดข้อความ PLO โดยอาศัยหลัก \textbf{Bloom's Revised Taxonomy} ควบคู่กับการเขียนวัตถุประสงค์ในรูปแบบ ABCD (Action verb, Behavior, Condition, Degree) เพื่อให้คำกริยาและผลลัพธ์มีความชัดเจนและวัดผลได้ สอดคล้องกับหลักปฏิบัติด้านประกันคุณภาพการศึกษาระหว่างประเทศและในประเทศ โดยเฉพาะอย่างยิ่ง การจัดความสัมพันธ์ของ PLO ต่อข้อความทิศทางของมหาวิทยาลัย เช่น วิสัยทัศน์ ``มหาวิทยาลัยพันธกิจสัมพันธ์ที่มีคุณภาพ สร้างคุณค่าเพื่อพัฒนาท้องถิ่น'' และปรัชญา IEL (Integrated Experiential Learning) ได้สรุปเป็นตารางที่ \ref{tbl:plo_vision_alignment} ส่วนระดับ Bloom ของแต่ละ PLO สรุปได้ที่ตารางที่ \ref{tbl:bloom_plo}

\textbf{ขั้นตอนการพัฒนา PLO ภายในคณะกรรมการพัฒนาหลักสูตร (สรุป)} มีลำดับดังนี้
\begin{criterionsubenumerate}
  \item ศึกษากรอบ AUN-QA version~4.0 และมาตรฐานคุณวุฒิระดับบัณฑิตศึกษา (TQF)
  \item วิเคราะห์วิสัยทัศน์/พันธกิจของมหาวิทยาลัยและคณะ รวมทั้งความต้องการของพื้นที่และผู้มีส่วนได้ส่วนเสียภายนอก
  \item ยกร่าง PLO และ PEO ตาม Bloom's Revised Taxonomy และรูปแบบ ABCD
  \item ผ่านการทบทวนโดยคณะกรรมการประจำคณะและสภาวิชาการ ก่อนเสนอสภามหาวิทยาลัยในมติครั้งที่ 4/2568
\end{criterionsubenumerate}

\begin{table}[H]
  \centering
  \caption{การสะท้อนเชิงธีมของ PLO ต่อวิสัยทัศน์ พันธกิจ/ภารกิจด้านผลผลิต และปรัชญา IEL (สรุปจากการออกแบบในมคอ.2)}
  \label{tbl:plo_vision_alignment}
  \begingroup
  \small
  \setlength{\tabcolsep}{3pt}
  \renewcommand{\arraystretch}{1.2}
  \begin{tabularx}{\linewidth}{|>{\raggedright\arraybackslash}X|c|c|c|c|c|}
    \hline
    \textbf{ข้อความทิศทาง (สรุป)} & \textbf{PLO1} & \textbf{PLO2} & \textbf{PLO3} & \textbf{PLO4} & \textbf{PLO5} \\
    \hline
    วิสัยทัศน์ มร.อ.: คุณภาพ พันธกิจสัมพันธ์ พัฒนาท้องถิ่น & \checkmark & \checkmark & \checkmark & \checkmark & \checkmark \\
    \hline
    พันธกิจ/ความจำเป็น: กำลังคนดิจิทัล--ปัญญาประดิษฐ์ และการแก้ปัญหาในพื้นที่ & \checkmark & \checkmark & \checkmark &  & \checkmark \\
    \hline
    ปรัชญา IEL: ประสบการณ์เชิงบูรณาการ & \checkmark & \checkmark & \checkmark & \checkmark & \checkmark \\
    \hline
  \end{tabularx}
  \endgroup
\end{table}

\noindent\textbf{ระดับ Bloom และการเลือกคำบรรยายผลการเรียนรู้ของแต่ละ PLO} ให้เป็นไปตามตารางที่ \ref{tbl:bloom_plo}

\begin{table}[H]
  \centering
  \caption{Bloom's Level ของแต่ละ PLO}
  \label{tbl:bloom_plo}
  \begin{tabularx}{\linewidth}{|>{\raggedright\arraybackslash}X|>{\raggedright\arraybackslash}X|>{\raggedright\arraybackslash}X|>{\raggedright\arraybackslash}X|}
    \hline
    \textbf{PLO} & \textbf{ชื่อย่อ} & \textbf{Bloom's Level} & \textbf{ระดับความซับซ้อน} \\
    \hline
    PLO1 & Create           & Create (C6)   & สร้างสรรค์สูงสุด \\
    PLO2 & Apply            & Apply (C3)    & ประยุกต์ใช้ \\
    PLO3 & Create Knowledge & Create (C6)   & สร้างสรรค์สูงสุด \\
    PLO4 & Ethics           & Evaluate (C5) & ประเมินค่า \\
    PLO5 & Analyze/Comm.    & Analyze (C4)  & วิเคราะห์ \\
    \hline
  \end{tabularx}
\end{table}

\noindent\textbf{ข้อความผลการเรียนรู้ที่คาดหวัง (PLOs) ของหลักสูตร}

\textbf{แผนที่ 1 (วิทยานิพนธ์) -- PLO1--PLO3 (Subject-specific)}
\begin{enumerate}
  \item \textbf{PLO1 (Create)} ใช้ซอฟต์แวร์หรือฮาร์ดแวร์ทางคอมพิวเตอร์ในการสร้างสรรค์ผลงาน องค์ความรู้ และนวัตกรรม
  \item \textbf{PLO2 (Apply)} ประยุกต์ใช้ความรู้ในสาขาวิชาวิศวกรรมคอมพิวเตอร์เพื่อแก้ไขปัญหาความต้องการขององค์กร ชุมชน และท้องถิ่น
  \item \textbf{PLO3 (Create Knowledge)} สร้างองค์ความรู้ทางด้านวิศวกรรมคอมพิวเตอร์และปัญญาประดิษฐ์ที่ตอบสนองความต้องการขององค์กร
\end{enumerate}

\textbf{แผนที่ 2 (สารนิพนธ์) -- PLO1--PLO5 (Subject-specific + Generic)}
\begin{enumerate}
  \item \textbf{PLO1 (Create)} ใช้ซอฟต์แวร์หรือฮาร์ดแวร์ทางคอมพิวเตอร์ในการสร้างสรรค์ผลงาน องค์ความรู้ และนวัตกรรม
  \item \textbf{PLO2 (Apply)} ประยุกต์ใช้ความรู้ในสาขาวิชาวิศวกรรมคอมพิวเตอร์เพื่อแก้ไขปัญหาความต้องการขององค์กร ชุมชน และท้องถิ่น
  \item \textbf{PLO3 (Create Knowledge)} สร้างต้นแบบนวัตกรรมที่ตอบสนองความต้องการขององค์กร
  \item \textbf{PLO4 (Ethics)} ปฏิบัติตามหลักจริยาบรรณวิชาชีพวิศวกรรมคอมพิวเตอร์
  \item \textbf{PLO5 (Analyze/Communicate/Collaborate/Lifelong)} สามารถคิดวิเคราะห์ การสื่อสาร การทำงานร่วมกัน การปรับตัว และการเรียนรู้ตลอดชีวิต
\end{enumerate}

ความสอดคล้องระหว่าง PEO และ PLO ในแผนที่ 2 แสดงให้เห็นว่าเป้าระดับโปรแกรมได้ถูกแตกย่อยลงเป็น PLO ที่วัดได้ ดังตารางที่ \ref{tbl:peo_plo_plan2}

\begin{table}[H]
  \centering
  \caption{ความเชื่อมโยง PEO--PLO -- แผนที่ 2 (สารนิพนธ์)}
  \label{tbl:peo_plo_plan2}
  \begin{tabularx}{\linewidth}{|>{\raggedright\arraybackslash}X|c|c|c|c|c|}
    \hline
    \rowcolor{gray!15}
    \textbf{PEO} & \textbf{PLO1} & \textbf{PLO2} & \textbf{PLO3} & \textbf{PLO4} & \textbf{PLO5} \\
    \hline
    PEO1 (ทักษะซอฟต์แวร์/นวัตกรรม) & \checkmark &            &            &            &            \\
    PEO2 (พัฒนาองค์ความรู้/ต้นแบบ) &            & \checkmark & \checkmark &            &            \\
    PEO3 (จริยธรรมวิชาชีพ)         &            &            &            & \checkmark &            \\
    PEO4 (เรียนรู้ตลอดชีวิต)       &            &            &            &            & \checkmark \\
    \hline
  \end{tabularx}
\end{table}

\noindent\textbf{การสื่อสาร PLO เพิ่มเติมหลังเริ่มรับเข้าศึกษา} เมื่อมีนักศึกษารุ่นแรก หลักสูตรใช้กิจกรรมปฐมนิเทศนักศึกษาระดับบัณฑิตศึกษา ปีการศึกษา 2568 เป็นช่องทางหลักในการถ่ายทอดข้อมูลหลักสูตรและ PLO แก่นักศึกษาโดยตรง โดย\textbf{เอกสารปฐมนิเทศ} (อ้างอิง AUNQA-1-5) ครอบคลุมโครงสร้างหลักสูตร ข้อบังคับบัณฑิตศึกษา และแนวทางการศึกษา ทำหน้าที่แทนคู่มือนักศึกษา ร่วมกับช่องทางเผยแพร่ออนไลน์ข้างต้น (รายละเอียดการรับฟังความต้องการจากผู้มีส่วนได้ส่วนเสียภายนอก ดูเกณฑ์ย่อย 1.4)

\par\vspace{4pt}\noindent%
\fcolorbox{teal}{teal!10}{%
\parbox{\linewidth-2\fboxsep-2\fboxrule}{%
\noindent\textbf{[RESOLVED] หลักฐานประกอบ Criterion~1 — สถานะ ณ 17 พ.ค. 2568:}%
\medskip%
\begin{itemize}\itemsep2pt
\item[(i)] \textbf{AUNQA-1-1} (ไฟล์ \texttt{มคอ 2 หลักสูตร...ผ่านสภา 2...docx}) --- URL เว็บมหาวิทยาลัย/ภาพหน้าจอ: \textbf{อัปเดตภายหลัง} (พร้อมหลักฐานแล้วคือไฟล์ มคอ.2 ฉบับผ่านสภา)
\item[(i-bis)] เว็บไซต์หลักสูตร \href{https://syweb.kidcbc.work}{\texttt{syweb.kidcbc.work}} แสดงผลลัพธ์การเรียนรู้ที่คาดหวัง (PLOs) ทั้ง 5 ข้อ พร้อม Bloom's Level ครบถ้วน — \textbf{[RESOLVED]} หลักฐาน Screenshot: \texttt{AUNQA-1-1b\_plo\_section\_syweb.png} (19 พ.ค. 2569)
\item[(ii)] เล่ม/ไฟล์ PDF คู่มือหรือเอกสารแนะนำหลักสูตร: \textbf{[RESOLVED]} เอกสารปฐมนิเทศนักศึกษาระดับบัณฑิตศึกษา ปีการศึกษา 2568 ระบุชื่อหลักสูตรวิศวกรรมศาสตรมหาบัณฑิต สาขาวิชาคอมพิวเตอร์และปัญญาประดิษฐ์ ครอบคลุมโครงสร้างหลักสูตร ข้อบังคับบัณฑิตศึกษา และแนวทางการศึกษา — หลักฐาน: \texttt{AUNQA-1-5\_ปฐมนิเทศนักศึกษาบัณฑิตศึกษา\_2568.pdf}
\item[(iii)] เพจ Facebook \href{https://www.facebook.com/profile.php?id=61577236346855}{วิศวกรรมคอมพิวเตอร์และปัญญาประดิษฐ์}: \textbf{[RESOLVED]} มีผู้ติดตาม 491 คน, ข้อมูลติดต่อ, โพสต์ประชาสัมพันธ์หลักสูตร — หลักฐาน: \texttt{AUNQA-1-1d\_facebook\_page\_profile.png} (19 พ.ค. 2569)
\item[(iv)] รายงานการประชุมคณะกรรมการพัฒนาหลักสูตร (Bloom/ABCD): \textbf{อัปเดตภายหลัง} (AUNQA-1-2 ระบุเป็น "รายงานการประชุมคณะกรรมการประจำหลักสูตร")
\item[(v)] ผลสำรวจผู้ประกอบการและการสนทนากลุ่ม: \textbf{อัปเดตภายหลัง} (กำลังดำเนินการรวบรวมจากฝ่ายบริหารหลักสูตร)
\end{itemize}%
}%
}\par\vspace{4pt}

% ------------------------------------------------------------------------------
\criterionsub{1.2 The programme to show that the expected learning outcomes for all courses are appropriately formulated and are aligned to the expected learning outcomes of the programme.}

ผลการเรียนรู้ของหลักสูตร (PLOs) ได้ถูกถ่ายทอดกระจายไปสู่รายวิชาบังคับ 6 รายวิชา และวิทยานิพนธ์/สารนิพนธ์ ผ่านกระบวนการทำ CLO (Course Learning Outcomes) โดยอาจารย์ผู้รับผิดชอบรายวิชาแต่ละท่าน กำหนด CLO ตามหลัก ABCD และเชื่อมโยง CLO เข้ากับ PLO อย่างน้อย 1 ข้อ ผ่านตาราง CLO--PLO Mapping ที่บรรจุไว้ใน มคอ.2 และในเอกสาร ``รวมเล่ม CLO รายวิชา'' (AUNQA-1-3)

\begin{table}[H]
  \centering
  \caption{การกระจายการเรียนรู้ที่คาดหวังของหลักสูตรลงสู่รายวิชาบังคับ}
  \begingroup
  \fontbody
  \setlength{\tabcolsep}{5pt}
  \renewcommand{\arraystretch}{1.2}
  \begin{tabularx}{\linewidth}{|>{\raggedright\arraybackslash}X|c|c|c|c|c|}
    \hline
    \textbf{รายวิชา} & \textbf{PLO1} & \textbf{PLO2} & \textbf{PLO3} & \textbf{PLO4} & \textbf{PLO5} \\
    \hline
    4095101 ปัญญาประดิษฐ์และการเรียนรู้ของเครื่อง & \checkmark &            & \checkmark &            &            \\
    4095102 การเขียนโปรแกรมขั้นสูงสำหรับ ML       & \checkmark &            & \checkmark &            &            \\
    7015906 ระเบียบวิธีวิจัยฯ                      &            & \checkmark & \checkmark & \checkmark & \checkmark \\
    7015907 สัมมนาวิศวกรรม CE-AI                   &            & \checkmark & \checkmark & \checkmark & \checkmark \\
    7015101 เทคโนโลยีอุบัติใหม่ฯ                   & \checkmark & \checkmark &            &            &            \\
    7015908 หัวข้อพิเศษฯ                           & \checkmark & \checkmark & \checkmark & \checkmark &            \\
    7015901--3 วิทยานิพนธ์ 1--3 (แผนที่ 1)         & \checkmark & \checkmark & \checkmark & \checkmark & \checkmark \\
    7015904--5 สารนิพนธ์ 1--2 (แผนที่ 2)           & \checkmark & \checkmark & \checkmark & \checkmark & \checkmark \\
    \hline
  \end{tabularx}
  \endgroup
\end{table}

% ------------------------------------------------------------------------------
\criterionsub{1.3 The programme to show that the expected learning outcomes consist of both generic outcomes (related to written and oral communication, problem-solving, information technology, teambuilding skills, etc) and subject specific outcomes (related to knowledge and skills of the study discipline).}

ผลการเรียนรู้ที่คาดหวังของหลักสูตรประกอบด้วย Generic Learning Outcomes และ Subject Specific Learning Outcomes ดังตารางต่อไปนี้

\begin{table}[H]
  \centering
  \caption{Generic Learning Outcomes และ Subject Specific Learning Outcomes}
  \begin{tabularx}{\linewidth}{|c|Y|c|c|}
    \hline
    \rowcolor{gray!15}
    \textbf{PLO} & \textbf{คำอธิบายผลการเรียนรู้} & \textbf{ประเภท} & \textbf{Level} \\
    \hline
    PLO1 & ใช้ซอฟต์แวร์/ฮาร์ดแวร์สร้างสรรค์ผลงาน นวัตกรรม & Subject specific & C \\
    \hline
    PLO2 & ประยุกต์ความรู้แก้ปัญหาขององค์กร/ชุมชน/ท้องถิ่น & Subject specific & A \\
    \hline
    PLO3 & สร้างองค์ความรู้/ต้นแบบนวัตกรรม CE \& AI & Subject specific & C \\
    \hline
    PLO4 & ปฏิบัติตามจริยาบรรณวิชาชีพวิศวกรรมคอมพิวเตอร์ & Generic & E \\
    \hline
    PLO5 & คิดวิเคราะห์ สื่อสาร ทำงานร่วมกัน เรียนรู้ตลอดชีวิต & Generic & AN \\
    \hline
  \end{tabularx}
\end{table}

\noindent\textbf{หมายเหตุ Level learning:} R = Remembering, U = Understanding, A = Applying, AN = Analyzing, E = Evaluating, C = Creating

% ------------------------------------------------------------------------------
\criterionsub{1.4 The programme to show that the requirements of the stakeholders, especially the external stakeholders, are gathered, and that these are reflected in the expected learning outcomes.}

หลักสูตรกำหนดผู้มีส่วนได้ส่วนเสียออกเป็น 4 กลุ่มหลัก โดยสามารถจัดเป็นสองมิติ ได้แก่ \textbf{ผู้มีส่วนได้ส่วนเสียภายใน} (อาจารย์ผู้รับผิดชอบหลักสูตร/อาจารย์ผู้สอน และนักศึกษาเมื่อเริ่มเปิดรับแล้ว) และ \textbf{ผู้มีส่วนได้ส่วนเสียภายนอก} (ผู้ใช้บัณฑิตในสถานประกอบการด้านดิจิทัล/AI ในภาคเหนือตอนล่าง องค์กรปกครองส่วนท้องถิ่นในจังหวัดอุตรดิตถ์และจังหวัดใกล้เคียง รวมถึงศิษย์เก่าหลักสูตรวิศวกรรมคอมพิวเตอร์ระดับปริญญาตรี) เพื่อให้การรวบรวมความต้องการครอบคลุมทั้งมุมมองผู้จัดการเรียนรู้ ผู้เรียน และผู้ใช้บัณฑิตในระบบเศรษฐกิจและสังคม

วิธีการได้มาซึ่งความต้องการของผู้มีส่วนได้ส่วนเสียประกอบด้วย (1)~\textbf{แบบสำรวจออนไลน์} ผ่าน Google Forms (อ้างอิง AUNQA-1-4b) รวบรวมความคิดเห็นจากกลุ่มผู้สนใจศึกษาต่อจำนวน \textbf{31 คน} แบ่งเป็นผู้กำลังศึกษาอยู่ร้อยละ 54.8 และผู้ที่ทำงานแล้วร้อยละ 45.2 โดยในจำนวนนี้มีผู้มีส่วนได้ส่วนเสียภายนอกจากหน่วยงานต่าง ๆ ได้แก่ การไฟฟ้าฝ่ายผลิตแห่งประเทศไทย บริษัท Fujitsu Thailand มหาวิทยาลัยเทคโนโลยีราชมงคลล้านนา โรงพยาบาลอุตรดิตถ์ มหาวิทยาลัยสุโขทัยธรรมาธิราช และ ก.ศ.น. จังหวัดอุตรดิตถ์ (2)~การสนทนากลุ่ม (focus group) กับศิษย์เก่าและผู้ใช้บัณฑิตระหว่างยกร่างหลักสูตร (3)~การวิเคราะห์แนวโน้มอุตสาหกรรมจากเอกสาร AI Workforce Development Plan ของภาครัฐ หลังจากนั้นนำข้อมูลมาวิเคราะห์/จัดกลุ่มความต้องการ และสะท้อนเข้าสู่ PLOs ดังตารางที่~\ref{tbl:needs_plo}

ผลสำรวจพบว่าเป้าหมายหลักของผู้ตอบแบบสำรวจ ได้แก่ การต่อยอดองค์ความรู้ด้านวิศวกรรมคอมพิวเตอร์และปัญญาประดิษฐ์ (ร้อยละ 51.6) การเตรียมพร้อมเพื่อประกอบอาชีพ (ร้อยละ 45.2) และการแก้ปัญหาหรือพัฒนางานที่ทำอยู่ (ร้อยละ 22.6) นอกจากนี้ผู้ตอบส่วนใหญ่ถึงร้อยละ 83.9 ต้องการเรียนในช่วงวันหยุดสุดสัปดาห์ ซึ่งสะท้อนให้เห็นความต้องการหลักสูตรที่ยืดหยุ่นเหมาะกับคนทำงาน ตารางที่~\ref{tbl:needs_plo} แสดงความสัมพันธ์ระหว่างความต้องการของผู้มีส่วนได้ส่วนเสียและ PLOs ที่สะท้อนความต้องการดังกล่าว

\begin{table}[H]
  \centering
  \caption{Needs-to-PLO Traceability Matrix — ความต้องการ SHs สะท้อนสู่ PLOs}
  \label{tbl:needs_plo}
  \begingroup
  \small
  \setlength{\tabcolsep}{3pt}
  \renewcommand{\arraystretch}{1.2}
  \begin{tabularx}{\linewidth}{|>{\raggedright\arraybackslash}X|>{\raggedright\arraybackslash}p{3.5cm}|c|c|c|c|c|c|}
    \hline
    \rowcolor{gray!15}
    \textbf{ความต้องการ (Needs)} & \textbf{กลุ่ม SHs} & \textbf{\%} & \textbf{PLO1} & \textbf{PLO2} & \textbf{PLO3} & \textbf{PLO4} & \textbf{PLO5} \\
    \hline
    ต่อยอดองค์ความรู้ด้าน CE\&AI (ผลสำรวจ $n=31$) & ผู้สนใจเรียนต่อ + External SHs & 51.6 & \checkmark & & \checkmark & & \\
    \hline
    ใช้ในการประกอบอาชีพในอนาคต (ผลสำรวจ) & นักศึกษา/ผู้ทำงาน & 45.2 & & & & \checkmark & \checkmark \\
    \hline
    แก้ปัญหาหรือพัฒนางานที่ทำอยู่ (ผลสำรวจ) & กลุ่มทำงาน + External SHs & 22.6 & & \checkmark & & & \\
    \hline
    ทักษะพัฒนาระบบ AI/Software สำหรับอุตสาหกรรม & ผู้ประกอบการ (กฟผ., Fujitsu) & — & \checkmark & & \checkmark & & \\
    \hline
    ความสามารถแก้ปัญหาองค์กร/ชุมชนท้องถิ่น & อปท. + ผู้ประกอบการ & — & & \checkmark & & & \\
    \hline
    จริยธรรมวิชาชีพและความรับผิดชอบ & ผู้ประกอบการ + อาจารย์ & — & & & & \checkmark & \\
    \hline
    การสื่อสาร ทำงานเป็นทีม เรียนรู้ตลอดชีวิต & ผู้ประกอบการ + อาจารย์ & — & & & & & \checkmark \\
    \hline
  \end{tabularx}
  \endgroup
\end{table}

\marknote{R1.4 — หลักฐานที่ยังขาด: (1)~\textbf{[บางส่วน OK]} แบบสำรวจออนไลน์ 31 คน มีผู้ประกอบการ/External SHs 8+ คน แต่ยังขาด employer survey เฉพาะกลุ่ม (ผู้จ้างงาน/นายจ้างโดยตรง) อย่างน้อย $n\geq10$ ราย (2)~\textbf{[ยังขาด]} บันทึกการประชุมหรือสรุปผล focus group ระหว่างยกร่างหลักสูตร (AUNQA-1-2) — ขอจากประธานหลักสูตร (3)~\textbf{[ยังขาด]} รายงานการวิเคราะห์ความต้องการตลาดแรงงานด้าน AI ระดับภูมิภาค}

% ------------------------------------------------------------------------------
\criterionsub{1.5 The programme to show that the expected learning outcomes are achieved by the students by the time they gradua.}

หลักสูตรเป็นหลักสูตรใหม่ พ.ศ. 2568 เริ่มมีนักศึกษาเข้าศึกษา (ปีการศึกษา 2568 จำนวน 4 คน ภาคการศึกษาที่ 2) จึงยังไม่มีผลการบรรลุ PLOs ของผู้สำเร็จการศึกษาระดับหลักสูตร อย่างไรก็ตาม หลักสูตรได้ออกแบบระบบการประเมินการบรรลุ PLOs ครอบคลุม 4 มิติ ตามกรอบ AUN-QA v4.0 ได้แก่ \textbf{When, Method, Tool} และ \textbf{Person} ดังตารางที่~\ref{tbl:plo_assessment_plan}

\begin{table}[H]
  \centering
  \caption{แผนการประเมินการบรรลุ PLOs — กรอบ 4 มิติ (When/Method/Tool/Person)}
  \label{tbl:plo_assessment_plan}
  \begingroup
  \small
  \setlength{\tabcolsep}{3pt}
  \renewcommand{\arraystretch}{1.3}
  \begin{tabularx}{\linewidth}{|c|>{\raggedright\arraybackslash}X|>{\raggedright\arraybackslash}X|>{\raggedright\arraybackslash}X|>{\raggedright\arraybackslash}X|}
    \hline
    \rowcolor{gray!15}
    \textbf{PLO} & \textbf{When (เวลา)} & \textbf{Method (วิธีการ)} & \textbf{Tool (เครื่องมือ)} & \textbf{Person (ผู้ประเมิน)} \\
    \hline
    PLO1 & ก่อนสำเร็จการศึกษา & ประเมินวิทยานิพนธ์/สารนิพนธ์ & เกณฑ์คณะกรรมการสอบ & คณะกรรมการสอบ \\
    \hline
    PLO2 & ปลายแต่ละภาคการศึกษา & CLO Assessment รายวิชา & แบบประเมิน CLO & อาจารย์ผู้สอน \\
    \hline
    PLO3 & ก่อนสำเร็จการศึกษา & ประเมินวิทยานิพนธ์/สารนิพนธ์ & Rubric นวัตกรรม/งานวิจัย & คณะกรรมการสอบ \\
    \hline
    PLO4 & ก่อนสำเร็จการศึกษา & แบบประเมินตนเอง + ประเมินโดยอาจารย์ & แบบประเมินจริยธรรม (PLO4 rubric) & อาจารย์ที่ปรึกษา \\
    \hline
    PLO5 & ปลายแต่ละปีการศึกษา & แบบประเมินตนเอง + แบบสำรวจ & แบบประเมินทักษะ Generic (PLO5) & อาจารย์ที่ปรึกษา \\
    \hline
  \end{tabularx}
  \endgroup
\end{table}

ผลการประเมิน CLO รายวิชาจะถูกรวบรวมในทุกปลายภาคการศึกษา เพื่อจัดทำ \textbf{PLO Dashboard} สำหรับติดตามแนวโน้มการบรรลุ PLOs ระหว่างทาง ก่อนที่จะมีบัณฑิตสำเร็จการศึกษา นอกจากนี้ หลักสูตรวางแผนสำรวจบัณฑิตและผู้ใช้บัณฑิต 1 ปีหลังสำเร็จการศึกษา เพื่อยืนยันการบรรลุ PLOs ในบริบทการทำงานจริง

\textbf{ข้อจำกัด:} ยังไม่มีผลการบรรลุ PLOs ที่แสดงได้จริงเนื่องจากยังไม่มีบัณฑิต -- คะแนนได้รับจากความสมบูรณ์ของการออกแบบระบบ ไม่ใช่จากผลลัพธ์ที่วัดได้จริง หลักสูตรจะเริ่มรายงาน \textbf{CLO Attainment รายวิชา} ตั้งแต่ปลายภาคการศึกษาที่ 2/2568 เป็นต้นไป

\begin{selfassessmentbox}{2}{หลักสูตรกำหนด PLOs ครบถ้วนตาม Bloom's Taxonomy เชื่อมโยงกับวิสัยทัศน์/พันธกิจ และกระจายลงสู่รายวิชาผ่าน CLO--PLO Mapping แต่ยังไม่มีหลักฐานเชิงประจักษ์การบรรลุ PLOs ของผู้สำเร็จการศึกษา และข้อมูล stakeholder feedback เชิงปริมาณยังไม่ครบถ้วน — มีนักศึกษาเรียนอยู่แล้วจึงสามารถเริ่มเก็บผล CLO และความพึงพอใจได้ต่อเนื่อง}
\end{selfassessmentbox}

% ------------------------------------------------------------------------------
\subsection*{รายการหลักฐาน}

หลักฐานประกอบ Criterion~1 ประกอบด้วย AUNQA-1-1 เล่ม มคอ.2 หลักสูตรวิศวกรรมคอมพิวเตอร์และปัญญาประดิษฐ์ ฉบับผ่านสภา พร้อมภาพหน้าจอเว็บไซต์หลักสูตร (\texttt{syweb.kidcbc.work}) และเพจ Facebook หลักสูตร, AUNQA-1-2 รายงานการประชุมคณะกรรมการประจำหลักสูตร, AUNQA-1-3 รวมเล่ม CLO รายวิชา พร้อม CLO--PLO Mapping, AUNQA-1-4 มติสภามหาวิทยาลัยราชภัฏอุตรดิตถ์ ครั้งที่ 4/2568 และผลการสำรวจความต้องการของผู้มีส่วนได้ส่วนเสีย 31 คน และ AUNQA-1-5 เอกสารปฐมนิเทศนักศึกษาระดับบัณฑิตศึกษา ปีการศึกษา 2568 (ใช้แทนคู่มือนักศึกษา)
